\chapter{线程}

\section{线程定义}
\concept{为什么引入线程}
\begin{points}
  \item 引入进程是为了描述和实现多道程序并发执行。
  \item 引入线程是为了减少并发执行的时空开销，让同一应用内部的多个任务能更轻量地并发。
  \item 早期进程同时承担“资源拥有单位”和“调度分派单位”两个角色，创建和切换开销较大。
  \item 线程把二者分开：进程更多作为资源拥有单位，线程作为 CPU 调度的基本执行单位。
\end{points}
\plain{线程是进程里面更轻的执行流，同进程线程共享资源，但各自有栈、寄存器和 PC。}
\exam{常考进程/线程区别、用户级线程和内核级线程优缺点、多对一/一对一/多对多模型。}


\concept{线程的定义}
\begin{points}
  \item 线程是进程内的一个执行单元，是进程内的可调度实体。
  \item 一个线程由线程 ID、程序计数器、寄存器集合和栈组成。
  \item 同一进程内的线程共享代码段、数据段、打开文件、地址空间等资源。
  \item 线程是“轻量级进程”，但不是独立资源拥有者。
\end{points}
\plain{线程是进程里面更轻的执行流，同进程线程共享资源，但各自有栈、寄存器和 PC。}
\exam{常考进程/线程区别、用户级线程和内核级线程优缺点、多对一/一对一/多对多模型。}


\concept{进程与线程区别}
\begin{points}
  \item 进程是资源分配的基本单位，线程是 CPU 调度的基本单位。
  \item 同一进程内线程共享地址空间，线程之间通信方便但相互影响更大。
  \item 不同进程地址空间隔离，通信开销较大但保护性更好。
  \item 线程创建、撤销、切换通常比进程轻量，因为不需要切换完整地址空间和大量资源。
\end{points}
\plain{程序是静态代码，进程是这段代码跑起来后的动态活动，还带资源和现场。}
\exam{常考程序与进程区别、进程实体由程序段/数据段/PCB 构成、进程是资源分配基本单位。}


\concept{多线程优点}
\begin{points}
  \item 响应性更好：例如界面线程响应用户，后台线程执行耗时任务。
  \item 资源共享方便：同一进程内线程天然共享内存和文件。
  \item 经济性更好：创建和切换开销较小。
  \item 可伸缩性更好：多核系统中多个线程可并行执行。
\end{points}
\plain{线程是进程里面更轻的执行流，同进程线程共享资源，但各自有栈、寄存器和 PC。}
\exam{常考进程/线程区别、用户级线程和内核级线程优缺点、多对一/一对一/多对多模型。}


\section{线程实现}
\concept{用户级线程}
\begin{points}
  \item 用户级线程由用户态线程库管理，内核不知道线程存在。
  \item 优点：创建和切换无需进入内核，速度快；调度策略可由应用定制。
  \item 缺点：一个线程执行阻塞系统调用，可能导致整个进程阻塞；无法真正利用多核并行。
\end{points}
\plain{线程是进程里面更轻的执行流，同进程线程共享资源，但各自有栈、寄存器和 PC。}
\exam{常考进程/线程区别、用户级线程和内核级线程优缺点、多对一/一对一/多对多模型。}


\concept{内核级线程}
\begin{points}
  \item 内核级线程由操作系统内核管理，调度器直接调度线程。
  \item 优点：一个线程阻塞不影响同进程其他线程；多线程可在多核上并行。
  \item 缺点：线程创建、切换、同步需要内核参与，开销比用户级线程大。
\end{points}
\plain{线程是进程里面更轻的执行流，同进程线程共享资源，但各自有栈、寄存器和 PC。}
\exam{常考进程/线程区别、用户级线程和内核级线程优缺点、多对一/一对一/多对多模型。}


\concept{多线程模型}
\begin{points}
  \item 多对一模型：多个用户线程映射到一个内核线程，管理轻量但阻塞问题明显。
  \item 一对一模型：每个用户线程映射到一个内核线程，并行能力强但内核线程数量受限。
  \item 多对多模型：多个用户线程映射到多个内核线程，兼顾灵活性与并行能力，但实现复杂。
\end{points}
\plain{线程是进程里面更轻的执行流，同进程线程共享资源，但各自有栈、寄存器和 PC。}
\exam{常考进程/线程区别、用户级线程和内核级线程优缺点、多对一/一对一/多对多模型。}


\concept{线程库}
\begin{points}
  \item 线程库向程序员提供创建、管理、同步线程的 API。
  \item 常见线程库包括 POSIX Pthreads、Windows threads、Java threads。
  \item 线程库可以完全在用户态实现，也可以作为内核级线程接口的包装。
\end{points}
\plain{线程是进程里面更轻的执行流，同进程线程共享资源，但各自有栈、寄存器和 PC。}
\exam{常考进程/线程区别、用户级线程和内核级线程优缺点、多对一/一对一/多对多模型。}


\concept{可重入代码}
\begin{points}
  \item 可重入代码指在被多个执行流同时调用或被中断后再次进入时，仍能得到正确结果的代码。
  \item 它通常不依赖不可保护的全局变量，不在函数内部保留共享静态状态，或对共享状态进行同步保护。
  \item 多线程、信号处理和中断处理程序常要求代码具备可重入性。
  \item 直觉：可重入代码像“每个人带自己的草稿纸”，不会抢同一张草稿纸写中间结果。
\end{points}
\plain{可重入代码像“公共只读函数”，多个线程同时进来也不会互相污染。}
\exam{常考可重入代码特征：不修改自身代码、不依赖全局可变状态或对其加保护。}


